% GENERATED FILE. Edit canonical lessons, then run npm run generate:books.
% canonical-source-hash: fnv1a64:112858d5581d861e
% canonical-lessons: TE-C187-tenetiga, TE-C187-balli, TE-C187-purugu, TE-C187-piccuka, TE-C187-kokila

\chapter{Bee, Lizard, Insect, Sparrow, Koel}
\label{ch:te-bee-lizard-insect-sparrow-koel}
\hlchaptermodality{187}

\begin{chapteropening}
\textbf{By the end of this chapter:} \emph{I can say a bee, a lizard, an insect, a sparrow and a koel.}

Complete the last lesson of chapter 187: I can say a bee, a lizard, an insect, a sparrow and a koel.
\end{chapteropening}

\section[\texorpdfstring{\te{తేనెటీగ} (tēneṭīga)}{తేనెటీగ (tēneṭīga)}]{\te{తేనెటీగ} (tēneṭīga) — a bee}

\label{lesson:TE-C187-tenetiga}



\begin{quote}
Before the new one: say the Telugu for twins, then the Telugu for a baby.
\end{quote}

\subsection*{You'll want to know: \te{తేనెటీగ}}
\textbf{\te{తేనెటీగ}} — \emph{tēneṭīga} — \textquotedblleft{}a bee\textquotedblright{}.

\begin{grammarlens}[title={What you've built}]
One more word for this chapter.
\end{grammarlens}

\subsection*{Your turn}
\begin{itemize}
\raggedright
  \item \emph{Say it:} \emph{tēneṭīga}
  \item \emph{Say it:} \emph{tēneṭīga}, once more
  \item \emph{Say it:} say \emph{tēneṭīga}
  \item \emph{Recall:} say the Telugu for twins, then the Telugu for a baby, then say \emph{tēneṭīga} again
\end{itemize}

\begin{tcolorbox}[breakable,colback=teal!4,colframe=teal!35!black,title={Before you move on}]
What does \te{తేనెటీగ} mean? (\textquotedblleft{}A bee\textquotedblright{}.) Say it once more.
\end{tcolorbox}
\section[\texorpdfstring{\te{బల్లి} (balli)}{బల్లి (balli)}]{\te{బల్లి} (balli) — a lizard, a gecko}

\label{lesson:TE-C187-balli}



\begin{quote}
Before the new one: say the Telugu for a baby, then the Telugu for a bee.
\end{quote}

\subsection*{You'll want to know: \te{బల్లి}}
\textbf{\te{బల్లి}} — \emph{balli} — \textquotedblleft{}a lizard, a gecko\textquotedblright{}.

\begin{grammarlens}[title={What you've built}]
One more word for this chapter.
\end{grammarlens}

\subsection*{Your turn}
\begin{itemize}
\raggedright
  \item \emph{Say it:} \emph{balli}
  \item \emph{Say it:} \emph{balli}, once more
  \item \emph{Say it:} say \emph{balli}
  \item \emph{Recall:} say the Telugu for a baby, then the Telugu for a bee, then say \emph{balli} again
\end{itemize}

\begin{tcolorbox}[breakable,colback=teal!4,colframe=teal!35!black,title={Before you move on}]
What does \te{బల్లి} mean? (\textquotedblleft{}A lizard, a gecko\textquotedblright{}.) Say it once more.
\end{tcolorbox}
\section[\texorpdfstring{\te{పురుగు} (purugu)}{పురుగు (purugu)}]{\te{పురుగు} (purugu) — an insect, a worm}

\label{lesson:TE-C187-purugu}



\begin{quote}
Before the new one: say the Telugu for a bee, then the Telugu for a lizard.
\end{quote}

\subsection*{You'll want to know: \te{పురుగు}}
\textbf{\te{పురుగు}} — \emph{purugu} — \textquotedblleft{}an insect, a worm\textquotedblright{}.

\begin{grammarlens}[title={What you've built}]
One more word for this chapter.
\end{grammarlens}

\subsection*{Your turn}
\begin{itemize}
\raggedright
  \item \emph{Say it:} \emph{purugu}
  \item \emph{Say it:} \emph{purugu}, once more
  \item \emph{Say it:} say \emph{purugu}
  \item \emph{Recall:} say the Telugu for a bee, then the Telugu for a lizard, then say \emph{purugu} again
\end{itemize}

\begin{tcolorbox}[breakable,colback=teal!4,colframe=teal!35!black,title={Before you move on}]
What does \te{పురుగు} mean? (\textquotedblleft{}An insect, a worm\textquotedblright{}.) Say it once more.
\end{tcolorbox}
\section[\texorpdfstring{\te{పిచ్చుక} (piccuka)}{పిచ్చుక (piccuka)}]{\te{పిచ్చుక} (piccuka) — a sparrow}

\label{lesson:TE-C187-piccuka}



\begin{quote}
Before the new one: say the Telugu for a lizard, then the Telugu for an insect.
\end{quote}

\subsection*{You'll want to know: \te{పిచ్చుక}}
\textbf{\te{పిచ్చుక}} — \emph{piccuka} — \textquotedblleft{}a sparrow\textquotedblright{}.

\begin{grammarlens}[title={What you've built}]
One more word for this chapter.
\end{grammarlens}

\subsection*{Your turn}
\begin{itemize}
\raggedright
  \item \emph{Say it:} \emph{piccuka}
  \item \emph{Say it:} \emph{piccuka}, once more
  \item \emph{Say it:} say \emph{piccuka}
  \item \emph{Recall:} say the Telugu for a lizard, then the Telugu for an insect, then say \emph{piccuka} again
\end{itemize}

\begin{tcolorbox}[breakable,colback=teal!4,colframe=teal!35!black,title={Before you move on}]
What does \te{పిచ్చుక} mean? (\textquotedblleft{}A sparrow\textquotedblright{}.) Say it once more.
\end{tcolorbox}
\section[\texorpdfstring{\te{కోకిల} (kōkila)}{కోకిల (kōkila)}]{\te{కోకిల} (kōkila) — a koel, the Indian cuckoo}

\label{lesson:TE-C187-kokila}



\begin{quote}
Before the new one: say the Telugu for an insect, then the Telugu for a sparrow.
\end{quote}

\subsection*{You'll want to know: \te{కోకిల}}
\textbf{\te{కోకిల}} — \emph{kōkila} — \textquotedblleft{}a koel, the Indian cuckoo\textquotedblright{}.

\begin{grammarlens}[title={What you've built}]
One more word for this chapter.
\end{grammarlens}

\subsection*{Your turn}
\begin{itemize}
\raggedright
  \item \emph{Say it:} \emph{kōkila}
  \item \emph{Say it:} \emph{kōkila}, once more
  \item \emph{Say it:} say \emph{kōkila}
  \item \emph{Recall:} say the Telugu for an insect, then the Telugu for a sparrow, then say \emph{kōkila} again
\end{itemize}

\begin{tcolorbox}[breakable,colback=teal!4,colframe=teal!35!black,title={Before you move on}]
What does \te{కోకిల} mean? (\textquotedblleft{}A koel, the Indian cuckoo\textquotedblright{}.) Say it once more.
\end{tcolorbox}
